\documentclass[12pt,a4paper]{article}

% ---------------------------------------------------------------
% Pacotes
% ---------------------------------------------------------------
\usepackage[utf8]{inputenc}
\usepackage[T1]{fontenc}
\usepackage[brazilian]{babel}
\usepackage{lmodern}
\usepackage[margin=2.5cm]{geometry}
\usepackage{amsmath,amssymb}
\usepackage{booktabs}
\usepackage{array}
\usepackage{graphicx}
\usepackage{listings}
\usepackage{xcolor}
\usepackage{caption}
\usepackage{enumitem}
\usepackage[hidelinks]{hyperref}

% ---------------------------------------------------------------
% Estilo dos trechos de código
% ---------------------------------------------------------------
\definecolor{codebg}{RGB}{245,245,245}
\definecolor{codekw}{RGB}{0,0,160}
\definecolor{codecm}{RGB}{90,120,90}
\lstset{
  language=Python,
  basicstyle=\ttfamily\small,
  keywordstyle=\color{codekw}\bfseries,
  commentstyle=\color{codecm}\itshape,
  backgroundcolor=\color{codebg},
  frame=single,
  framerule=0pt,
  breaklines=true,
  showstringspaces=false,
  columns=fullflexible,
  keepspaces=true,
  literate=
    {á}{{\'a}}1 {é}{{\'e}}1 {í}{{\'i}}1 {ó}{{\'o}}1 {ú}{{\'u}}1
    {â}{{\^a}}1 {ê}{{\^e}}1 {ô}{{\^o}}1 {ã}{{\~a}}1 {õ}{{\~o}}1
    {ç}{{\c{c}}}1 {à}{{\`a}}1
}

\newcommand{\cp}{\,\text{cp}}

% ---------------------------------------------------------------
% Documento
% ---------------------------------------------------------------
\begin{document}

% ===============================================================
% CAPA
% ===============================================================
\begin{titlepage}
  \centering
  \vspace*{1.5cm}

  {\Large\textsc{Instituto Tecnológico de Aeronáutica}\par}

  \vspace{4cm}

  {\Huge\bfseries Projeto Xadrez Heurístico\par}
  \vspace{0.5cm}
  {\Large Relatório do Laboratório\par}

  \vspace{4cm}

  {\large\bfseries Membros da Equipe\par}
  \vspace{0.4cm}
  {\large
    Gabriel Negreiro\par
    João Gabriel\par
    Marcello Ryaj\par
    Mateus Levi\par
  }

  \vfill

  {\large CMC-15 -- Inteligência Artificial\par}
  {\large Prof. Paulo André L. de Castro\par}
  \vspace{0.6cm}
  {\large São José dos Campos\par}
  {\large Setembro de 2026\par}
\end{titlepage}

% ===============================================================
% FUNÇÃO HEURÍSTICA
% ===============================================================
\section{Função Heurística}

A função heurística atribui um valor numérico a uma posição quando a
busca não alcança um estado terminal; o Minimax apenas propaga para a
raiz os valores que ela produz nas folhas. Adotou-se a convenção de que
\textbf{valores positivos favorecem as brancas} e \textbf{negativos
favorecem as pretas}. A função recebe uma posição e um lance candidato,
aplica o lance numa cópia do tabuleiro e avalia a posição resultante.

A maior parte dos termos e dos pesos foi retirada da \emph{Chess
Programming Wiki}, em especial da \emph{Simplified Evaluation Function}
de Michniewski~\cite{cpw-sef}, e das páginas sobre mobilidade e estrutura
de peões~\cite{cpw-mobility,cpw-pawn}. Aqui descrevemos apenas o
essencial de cada termo e as decisões que foram nossas; os detalhes das
tabelas e a justificativa de cada valor estão nas referências.

\subsection{Problemas da função original}

A versão fornecida tinha três defeitos que comprometiam a busca:
\textbf{(i)} todo xeque-mate valia $-100$, sem distinguir quem dava o
mate, então o agente nunca buscava o mate ativamente;
\textbf{(ii)} qualquer xeque retornava $-5$ imediatamente, descartando a
contagem de material;
\textbf{(iii)} o empate valia $-10$, pior que muitas posições perdidas.
Além disso, a avaliação posicional se resumia ao controle das quatro
casas centrais e a um termo simples de segurança do rei.

\subsection{Valores terminais}

Após aplicar o lance, \texttt{board.turn} indica o lado a jogar, que em
um xeque-mate é o perdedor. Assim, mate com as pretas a jogar vale
$+100\,000$ e com as brancas a jogar vale $-100\,000$; o valor domina o
máximo material possível ($\approx 3\,900\cp$). Todo outro fim de jogo
vale exatamente $0$, de modo que o agente prefere empate a derrota e
aceita o empate quando está em desvantagem. O xeque deixou de ser um
retorno antecipado e virou apenas uma penalidade dentro da segurança do
rei.

\subsection{Material e tabelas posicionais}

Os valores de material foram convertidos para centipawns conforme
\cite{cpw-sef}: peão 100, cavalo 320, bispo 330, torre 500, dama 900 e
rei 0 (ambos os lados sempre têm um). Da mesma fonte vieram as
\emph{piece-square tables}: um bônus por casa para cada tipo de peça,
que substitui o controle de centro original e cobre também
desenvolvimento e roque, já que casas iniciais pontuam mal e o rei
pontua bem em \texttt{g1}/\texttt{c1}. As tabelas são escritas do ponto
de vista das brancas com a oitava fileira primeiro; como a biblioteca
\texttt{chess} numera as casas a partir de \texttt{a1 = 0}, as brancas
indexam com \texttt{square\_mirror(square)} e as pretas com
\texttt{square} diretamente.

\subsection{Mobilidade}

Diferença entre o número de lances legais das brancas e das pretas,
ponderada por $5\cp$ por lance~\cite{cpw-mobility}. Como
\texttt{legal\_moves} só enumera o lado a jogar, a contagem é feita numa
cópia do tabuleiro invertendo \texttt{turn}; alterar o turno no
tabuleiro real corromperia a partida.

\subsection{Estrutura de peões}

A partir de uma contagem de peões por coluna, penalizam-se
\textbf{peões dobrados} ($n$ peões na mesma coluna contam $n-1$
defeitos) e \textbf{peões isolados} (sem peão amigo nas colunas
adjacentes, com as bordas \texttt{a} e \texttt{h} tratadas como vazias
para evitar que \texttt{structure[n-1]} dê a volta em Python). Cada
defeito custa $50\cp$, o $-0{,}5$ peão de Shannon~\cite{cpw-pawn}.

\subsection{Segurança das peças}

Uma peça está \emph{pendurada} quando é atacada e não tem defensor
(\texttt{attackers} da cor oposta $>0$ e da própria cor $=0$; o rei é
ignorado, pois rei atacado é xeque). A penalidade é de \textbf{10\% do
valor da peça}, não 100\%: dentro do horizonte o Minimax já enxerga a
captura real, e penalizar integralmente contaria a perda duas vezes. O
termo apenas sinaliza o risco de peças capturáveis além do horizonte.

\subsection{Composição final}

Defeitos são subtraídos para as brancas e somados para as pretas:
\begin{align*}
  f(s) ={}& \text{material} + \text{PST}
          + 5\,(\text{mob}_B - \text{mob}_P)
          - 50\,(\text{def}_B - \text{def}_P) \\
        & - 0{,}1\,(\text{pend}_B - \text{pend}_P)
          + (\text{rei}_B - \text{rei}_P),
\end{align*}
onde o termo de segurança do rei, mantido da versão original e
reescalado, dá $-15\cp$ por casa vazia no escudo de peões e $-50\cp$ se
o lado está em xeque.

\subsection{Testes}

Cada termo foi validado em posições controladas cujo valor esperado foi
conferido manualmente. Todos os casos passaram.

\begin{table}[h]
  \centering
  \caption{Casos de teste da função heurística.}
  \small
  \begin{tabular}{>{\raggedright\arraybackslash}p{3.4cm} >{\raggedright\arraybackslash}p{6.8cm} r}
    \toprule
    Caso & Posição & Esperado = Obtido \\
    \midrule
    Mate dado pelas brancas
      & Mate do pastor: \texttt{1.e4 e5 2.Bc4 Nc6 3.Qh5 Nf6 4.Qxf7\#}
      & $100\,000$ \\
    Mate dado pelas pretas
      & Mate do bobo: \texttt{1.f3 e5 2.g4 Qh4\#}
      & $-100\,000$ \\
    Afogamento
      & \texttt{7k/8/6K1/8/8/8/8/5Q2 w}, lance \texttt{Qf7}
      & $0$ \\
    Peões dobrados
      & Brancas com peões em a2, b2, b3, e2, h2
      & $1$ \\
    Peões isolados
      & Mesma posição: colunas \texttt{e} e \texttt{h} sem vizinhas
      & $2$ \\
    Penduradas das pretas
      & \texttt{r3k3/8/3p4/4N3/8/2b5/3Q4/4K3 w}: peão d6 e bispo c3, $0{,}1\times(100+330)$
      & $43{,}0$ \\
    Penduradas das brancas
      & Mesma posição: cavalo e5, $0{,}1\times 320$
      & $32{,}0$ \\
    \bottomrule
  \end{tabular}
\end{table}

% ===============================================================
% REFERÊNCIAS
% ===============================================================
\begin{thebibliography}{9}
\raggedright

\bibitem{cpw-sef}
Chess Programming Wiki. \emph{Simplified Evaluation Function}
(Tomasz Michniewski).
Disponível em: \url{https://www.chessprogramming.org/Simplified_Evaluation_Function}.
Acesso em: set.\ 2026.

\bibitem{cpw-mobility}
Chess Programming Wiki. \emph{Mobility}.
Disponível em: \url{https://www.chessprogramming.org/Mobility}.
Acesso em: set.\ 2026.

\bibitem{cpw-pawn}
Chess Programming Wiki. \emph{Pawn Structure}.
Disponível em: \url{https://www.chessprogramming.org/Pawn_Structure}.
Acesso em: set.\ 2026.

\end{thebibliography}

\end{document}
